\chapter{Conclusion and Next Steps}
\sffamily

In summary, the ATPW team within the MFWF project has completed several of the key objectives set out at the beginning of the semester. The aerodynamic test platform has a completed inner frame for mounting of the MFWF in such a way that transmits only the orthogonal forces to the load cells. Rigorous testing has been performed throughout the semester to verify the functionality of the aerodynamic test platform. While major design changes were needed in terms of the electronics, the platform is now capable of providing accurate force readings down to 30 grams of applied static load. In the event of a required redesign of the platform, preliminary designs have been developed in the attempts to rasing the resolution of the measured forces.

The ATPW team has achieved a fully functional and accurate MLP neural network able to predict the CL, CD, CP, and CDP of multiple NACA airfoils. The network has been trained on a large dataset of over 66,000 samples of NACA airfoil data, and has achieved training, validation, and test accuracies of 90.48\%, 91.88\%, and 96.30\%, respectively.

 A full redesign of the wing architecture was executed to meet the revised interface requirements of the electromagnetic flapping mechanism, a structured Design of Experiments spanning two planform series, three vein splay patterns, and three chord lengths was designed, and all 40 wings were manufactured, quality controlled, and dynamically screened through drill spin testing at 30--33~Hz. A repeatable fabrication process was validated. A provisional optimal wing, R-V-C5.0 (ID:~F), was identified based on mass compliance and qualitative camber performance, pending ATP validation. The dynamic testing and motion tracking framework also established. This included the development of a dual-camera imaging system and a MATLAB-based tracking pipeline capable of extracting both in-plane and out-of-plane deformation using colour-based marker tracking. While full testing of the wing set was not completed within the project timeline, the developed framework provides a practical foundation for future validation of wing designs.

In the Winter 2026 semester, the ATPW team will be starting off strong by focusing on three main objectives. The first will be to fully commission the aerodynamic test platform with a representative MFWF model. This will involve mounting a test flyer to the platform and performing a series of known dynamic loadings to verify the platform's ability to accurately measure forces during flapping.
The wings team will fabricate and test the new wing designs and continue building the dynamic performance database that will support the selection of the final functional wing pair. The neural network team will focus on adapting the MLP model to predict the three-dimensional aerodynamic forces acting on the flapping wings of the MFWF. This will require collaboration with the rest of the ATPW team to collect sufficient data from the test platform to train and validate the model effectively.

FUTURE WORK PARAGRAPH (Abdul and justin)
Future work in terms of wing design and manufacturing includes defining an interface with the concerned stakeholders, expanding the DOE or testing the existing wing set, exploring more membranes and completing ATP testing. The primary outcome of dynamic testing and motion tracking is the development of a functional dynamic testing and tracking framework, with future work focused on generating repeatable and comparable datasets across multiple wing designs. Key improvements include the use of high-speed imaging, optimized marker placement, and a more consistent flapping mechanism to improve tracking accuracy and data reliability.